\chapter{Software Requirement Specification}
\label{ch:srs}

This chapter specifies the requirements of the benchmark system: the corpus
construction pipeline, the mutator suite, and the evaluation harness. It states
the overall description of the system, its perspective and functions, and the
functional, performance, software, and hardware requirements that the
implementation satisfies. The chapter closes with a summary of the key
specifications.

%----------------------------------------------------------

\section{Overall Description}

The system is a research pipeline, not an interactive application. It ingests
raw vulnerability records, validates their labels, constructs safe
counterparts, generates adversarial variants, runs a set of detectors, and
emits scored results. Its users are security researchers and tool authors who
need a trustworthy substrate to measure Python vulnerability detectors.

\subsection{Product Perspective}

The benchmark is a standalone, reproducible pipeline with three stages that hand
off through versioned artifacts on disk. The first stage produces a validated
corpus of labeled samples. The second stage produces mutated variants of the
test samples. The third stage runs detectors and computes metrics. Each stage
reads and writes schema-versioned files, so a stage can be re-run in isolation
and any result can be traced back to the inputs that produced it. The system
relates to external systems only as data sources: advisory databases and CVE
feeds upstream, and detector binaries or APIs at evaluation time.

\subsection{Product Functions}

The system performs the following functions:

\begin{itemize}
  \item Scrape positive samples from five advisory and CVE-fix sources.
  \item Filter samples by sink presence, removing those without the
    category-defining sink token.
  \item Support an independent, CWE-stratified audit of label correctness.
  \item Apply a diff-based filter that keeps a positive only if its sink line
    changed between the vulnerable and fixed versions.
  \item Construct safe hard negatives as remediated twins of real vulnerable
    functions.
  \item Split the corpus into train, validation, and test sets with no
    cross-split repository leakage.
  \item Generate seven semantics-preserving mutated variants of each test
    sample.
  \item Run seven detectors across four tiers and score them under five bounded
    metrics on clean and mutated code.
\end{itemize}

\subsection{Constraints}

The system operates under these constraints:

\begin{itemize}
  \item Labels must be evidence-backed: a positive is admitted only with a
    matched sink pattern and, where a paired fix exists, a changed sink line.
  \item Every mutator must be semantics-preserving, and every mutated sample
    must re-parse as valid Python.
  \item Samples sharing a repository must fall in the same split, to prevent
    leakage.
  \item All validation decisions must be reproducible from released artifacts.
\end{itemize}

%----------------------------------------------------------

\section{Specific Requirements}

\subsection{Functional Requirements}

\begin{enumerate}
  \item The scraper shall retrieve per-CWE positives with explicit fix commits
    where the source provides them.
  \item The sink-presence filter shall admit a sample only if at least one
    regex sink pattern matches its source after comments and docstrings are
    stripped.
  \item The audit interface shall present each item with its identifier,
    source, repository, file path, matched pattern, and code excerpt, and
    record a binary verdict with a one-sentence justification.
  \item The diff filter shall compare the vulnerable and fixed versions after
    comment stripping and whitespace normalization, and reject a positive whose
    sink lines are unchanged.
  \item The negative builder shall apply the canonical remediation for a CWE at
    the AST level, preserving surface structure and the sink name while
    neutralizing the sink behavior.
  \item Each mutator shall transform the Python AST deterministically given a
    sample identifier and seed, and shall skip any node it cannot prove safe to
    transform.
  \item The harness shall score every detector one-versus-rest per CWE and
    report macro-F1, MCC, SWR, HCMA, and weighted $\kappa$.
\end{enumerate}

\subsection{Performance Requirements}

The pipeline runs offline in batch. The throughput requirement is that the full
corpus, on the order of two thousand pre-filter samples, be processed end to
end on a single workstation. Detector runs must complete within the rate limits
of the hosted LLM APIs, which sets the practical bound on wall-clock time. The
trained detector must fit and run on a single GPU. Determinism is required: the
SAST tools and the fixed-checkpoint model produce identical output across runs,
and the LLMs run at temperature zero to keep per-call drift small.

\subsection{Software Requirements}

\begin{itemize}
  \item Python 3.11 or later for the pipeline, mutators, and harness.
  \item The Python standard \texttt{ast} module for AST transformations.
  \item PyTorch and the HuggingFace Transformers library for the trained
    detector.
  \item Bandit 1.9.4, Semgrep 1.163.0, and Snyk Code 1.1305.0 as SAST and
    dataflow detectors.
  \item API access to Claude Sonnet 4.6, GPT-4o, and DeepSeek R1 for the LLM
    tier.
  \item Git and SQLite for source retrieval (CVEfixes ships as a SQLite
    database).
\end{itemize}

\subsection{Hardware Requirements}

\begin{itemize}
  \item A workstation with at least 16 GB of RAM for the pipeline and harness.
  \item A CUDA-capable GPU with at least 12 GB of memory to fine-tune the
    126.6M-parameter GraphCodeBERT detector.
  \item Sufficient disk for the corpus, mutated variants, and per-detector
    prediction logs.
\end{itemize}

\subsection{Design Constraints}

The single-file extract format is a deliberate design constraint: each sample is
a self-contained function, which suits SAST, LLM, and learned detectors but is
unfavorable to whole-project dataflow tools that need application entry points.
This trade-off is accepted so that the benchmark stays language-uniform and
reproducible, and the whole-project tools are retained in the harness for a
future application-level format.

\section{Summary}

The benchmark system is a three-stage, artifact-driven pipeline whose
requirements center on evidence-backed labels, semantics-preserving mutation,
leakage-free splits, and reproducible scoring. The functional requirements
define each pipeline step, and the non-functional requirements fix the software
and hardware needed to run it. These specifications guide the high-level design
in Chapter~\ref{ch:hld}.
